% !TeX root = ../../金融经济学_熊和平_习题.tex

\begin{problem}
\problemno{5.1} 什么是套利？套利一般有哪几个类别？
\end{problem}
\begin{solution}
套利是一个组合 $\theta$，满足初始成本 $S^{\mathsf T}\theta\le0$，且期末支付 $X\theta\ge0$、至少一个状态严格为正。按来源可分为：利用同一支付不同价格的空间套利，利用不同到期价格关系的时间套利，利用复制关系的证券套利，以及利用统计定价偏离的统计套利。前三类是无套利定价理论的直接对象。
\end{solution}
\begin{problem}
\problemno{5.2} 什么是无套利原理？无套利原理成立的前提假设是什么？
\end{problem}
\begin{solution}
无套利原理要求不存在满足上述不等式的组合。其常用前提是：证券可按任意实数数量线性组合，市场价格对投资者给定，卖空和借贷在模型中可行，交易成本、税收和流动性限制忽略，且投资者严格偏好每个状态都更多的财富。若存在交易摩擦，应改用带摩擦的无套利条件。
\end{solution}
\begin{problem}
\problemno{5.3} 如何利用无套利原理对证券进行定价？
\end{problem}
\begin{solution}
设目标证券支付为 $Y$。若存在组合 $\theta$ 使 $X\theta=Y$，则复制组合成本 $S^{\mathsf T}\theta$ 必须等于目标价格 $P$；否则买入便宜的一方、卖出昂贵的一方即可在支付相同的情况下获得正初始收入。若只能找到 $X\theta\ge Y$，则 $S^{\mathsf T}\theta$ 给出目标价格的上界；反向支配关系给出下界。
\end{solution}
\begin{problem}
\problemno{5.4} 为什么要讨论证券市场的不同描述？
\end{problem}
\begin{solution}
支付矩阵 $X$ 直接描述每支证券在每个状态的支付，价格向量 $S$ 描述期初价值，组合空间 $\{X\theta\}$ 描述可实现支付，状态价格 $q$ 则把支付映射为价格 $S=X^{\mathsf T}q$。支付矩阵适合判断秩和完全性，状态价格适合定价和风险中性转换，组合空间适合直接检验套利。
\end{solution}
\begin{problem}
\problemno{5.5} 回忆线性代数中的秩的概念，它在证券市场的描述中起何作用？
\end{problem}
\begin{solution}
支付矩阵的秩是列空间维数，即独立可交易支付的数量。若有 $S$ 个状态，$\operatorname{rank}(X)=S$ 时每个状态支付都可复制，市场完全；若秩小于 $S$，至少有一个状态支付不能由证券组合生成。若证券数 $N$ 大于秩，则存在非零 $\theta$ 使 $X\theta=0$，证券中有冗余组合。
\end{solution}
\begin{problem}
\problemno{5.6} 冗余证券在套利中起什么作用？
\end{problem}
\begin{solution}
若 $X_j=X_{-j}\alpha$，证券 $j$ 的支付可由其他证券组合复制。无套利要求 $S_j=S_{-j}^{\mathsf T}\alpha$。价格高于右侧时卖出证券 $j$、买入复制组合；价格低于右侧时反向交易。两种情况下期末净支付均为零，期初都有正收入，因此冗余证券是检验一价律的直接工具。
\end{solution}
\begin{problem}
\problemno{5.7} 什么是一价律？一价律在衍生证券定价中是如何应用的？
\end{problem}
\begin{solution}
一价律是指 $X\theta=X\theta'$ 必有 $S^{\mathsf T}\theta=S^{\mathsf T}\theta'$。对衍生证券，先求组合 $\theta$ 使 $X\theta=Y$，则其无套利价格为 $P=S^{\mathsf T}\theta$。例如看涨期权可以由股票和无风险证券复制；若只能构造支付处处不低于 $Y$ 的组合，则复制成本是价格上界。
\end{solution}
\begin{problem}
\problemno{5.8} 什么是资产定价基本定理？理解其“基本”的含义。
\end{problem}
\begin{solution}
在有限状态模型中，资产定价基本定理表明以下命题等价：市场无套利；存在严格正的状态价格向量 $q$ 使 $S=X^{\mathsf T}q$；存在与客观概率等价的风险中性概率，使所有资产的贴现期望支付等于当前价格。“基本”是因为它把套利、线性状态价格和风险中性定价三种表述统一起来，是后续衍生证券定价的逻辑基础。
\end{solution}

\begin{problem}
\problemno{5.9} 证明定理 5.8 中的 (1)。
\end{problem}
\begin{solution}
若存在严格正的 $q$ 满足 $S=X^{\mathsf T}q$，设 $\theta$ 是套利组合。由价格关系，
\[
S^{\mathsf T}\theta=(X^{\mathsf T}q)^{\mathsf T}\theta=q^{\mathsf T}X\theta.
\]
套利定义给出 $S^{\mathsf T}\theta\le0$，而 $X\theta\ge0$ 且至少一个分量严格为正。因为每个 $q_s>0$，所以 $q^{\mathsf T}X\theta>0$，从而 $S^{\mathsf T}\theta>0$，与套利的初始成本条件矛盾。因此不存在套利。
\end{solution}
\begin{problem}
\problemno{5.10} 证券市场上有三个证券，其支付矩阵为
\[
    X=\begin{bmatrix}
        30 & 20 & 10\\
        20 & 15 & 0
    \end{bmatrix}.
\]
\begin{enumerate}[label=(\alph*)]
    \item 是否存在冗余证券？哪个证券是冗余的？请证明之。
    \item 若存在套利机会，请给出套利策略。
\end{enumerate}
\end{problem}
\begin{solution}
\begin{enumerate}[label=(\alph*)]
\item 由于只有两个状态而有三个证券，三列必线性相关。具体地，
\[
X_3=3X_1-4X_2
=3\begin{bmatrix}30\\20\end{bmatrix}
-4\begin{bmatrix}20\\15\end{bmatrix}
=\begin{bmatrix}10\\0\end{bmatrix}.
\]
因此证券 3 是冗余证券；它的支付可以由买入 3 单位证券 1、卖空 4 单位证券 2 完全复制。
\item 设三证券价格为 $S_1,S_2,S_3$。复制组合的成本为 $3S_1-4S_2$。若
\[
S_3>3S_1-4S_2,
\]
则卖出 1 单位证券 3，同时买入 3 单位证券 1、卖空 4 单位证券 2，期初获得正现金 $S_3-3S_1+4S_2$，期末支付恒为零，构成套利。若 $S_3<3S_1-4S_2$，则反向交易构成套利。因此无套利的必要条件是 $S_3=3S_1-4S_2$；在该等式成立时，仅由这组支付关系不能再构造价格套利。
\end{enumerate}
\end{solution}
\begin{problem}
\problemno{5.11} 若市场上三个证券的支付矩阵为
\[
    X=\begin{bmatrix}
        110 & 100 & 48\\
        110 & 50 & 40\\
        110 & 40 & 36
    \end{bmatrix},
\]
\begin{enumerate}[label=(\alph*)]
    \item 市场是否完全？请证明之。
    \item 若三个证券的价格分别为 $100$、$70$ 和 $40$，则是否存在套利机会？如果不存在套利机会，请计算无风险利率和 A-D 证券价格以及等价鞅测度。
    \item 若三个证券的价格分别为 $100$、$40$ 和 $70$，则是否存在套利机会？
    \item 上述三个复合证券分别是由怎样的 A-D 证券组合复合而成的？
\end{enumerate}
\end{problem}
\begin{solution}
\begin{enumerate}[label=(\alph*)]
\item
\[
\det X=13200\ne0,
\]
所以三列线性无关，且秩等于状态数 $3$，市场完全。
\item 令 $q=(q_1,q_2,q_3)^{\mathsf T}$。由 $X^{\mathsf T}q=(100,70,40)^{\mathsf T}$ 得
\[
110(q_1+q_2+q_3)=100,\quad100q_1+50q_2+40q_3=70,
\]
\[
48q_1+40q_2+36q_3=40.
\]
解得 $q=(17/33,3/11,4/33)^{\mathsf T}>0$，故不存在套利。总状态价格为 $8/11$，所以
\[
1+r_f=\frac{11}{8},\qquad r_f=\frac38.
\]
等价鞅概率为 $\widehat p=q/(8/11)=(17/24,3/8,1/6)^{\mathsf T}$。
\item 对价格 $(100,40,70)$ 解同一方程会得到负的状态价格分量。由资产定价基本定理，不存在严格正状态价格，故存在套利；套利组合可由该线性系统的负系数方向构造。
\item 三个证券分别是三列支付向量对应的状态证券组合：
\[
X_1=110A_1+110A_2+110A_3,quad
X_2=100A_1+50A_2+40A_3,
\]
\[
X_3=48A_1+40A_2+36A_3.
\]
\end{enumerate}
\end{solution}
\begin{problem}
% 原书 p.108：支付矩阵 OCR 残损，按三状态三资产整理，待补扫核对。
\problemno{5.12} 已知某个市场的支付矩阵为
\[
    X=\begin{bmatrix}
        1 & 0 & 0\\
        1 & 2 & 1\\
        1 & 3 & b
    \end{bmatrix},
\]
相应的价格向量为 $S=(6,31,9)^{\mathrm{T}}$。
\begin{enumerate}[label=(\alph*)]
    \item 证明存在冗余证券；
    \item $b$ 为何值时不存在套利机会？
    \item 不存在套利机会时计算状态价格、风险中性概率和利率。
\end{enumerate}
\end{problem}
\begin{solution}
\begin{enumerate}[label=(\alph*)]
\item 三个证券的支付列向量位于三维状态空间；是否冗余要看列之间的线性关系。由矩阵可知 $X_1=(1,1,1)^{\mathsf T}$、$X_2=(0,2,3)^{\mathsf T}$、$X_3=(0,1,b)^{\mathsf T}$。当 $b=3/2$ 时 $X_3=\frac12X_2$，证券 3 冗余。
\item 由价格方程
\[
q_1+q_2+q_3=6,\quad 2q_2+3q_3=31,\quad q_2+bq_3=9,
\]
可得
\[
q_3=\frac{13}{3-2b},\quad q_2=\frac{27-31b}{3-2b},\quad q_1=\frac{19b-22}{3-2b}.
\]
严格正状态价格同时要求 $b<27/31$（由 $q_2>0$）和 $b>22/19$（由 $q_1>0$），两者不可能同时成立。因此按当前整理出的支付矩阵和价格，不存在任何 $b$ 能使市场无套利；这也说明原书此处的 OCR 矩阵必须进一步核对。
\item 总状态价格恒为 $6$，故 $1/(1+r_f)=6$，即 $r_f=-5/6$。风险中性概率为 $\widehat p_s=q_s/6$，将上式分别除以 6 即得。
\end{enumerate}
\end{solution}
\begin{problem}
\problemno{5.13} 给定支付矩阵
\[
    X=\begin{bmatrix}
        2 & 3 & 2\\
        2 & 3 & 5\\
        5 & 4/3 & 3
    \end{bmatrix},
\]
和资产价格 $S^{\mathrm{T}}=(6,4,k)$：
\begin{enumerate}[label=(\alph*)]
    \item $K$ 为何值时不存在套利机会？
    \item 取 $K$ 为某特定值使得市场存在三类套利机会，分别给出相应的套利组合。
    \item 求以第一个资产为标的、以 $3$ 为执行价格的欧式看涨期权价格；求以第三个资产为标的、以 $2.5$ 为执行价格的欧式看涨期权价格。
    \item 若四个状态发生的概率相等，计算四个状态对应的随机贴现因子。
\end{enumerate}
\end{problem}
\begin{solution}
\begin{enumerate}[label=(\alph*)]
\item 令状态价格为 $q$，无套利条件为 $X^{\mathsf T}q=(6,4,k)^{\mathsf T}$ 且 $q>0$。先由前两列方程解出 $q_1,q_2,q_3$ 的关系，再将第三列代入得到 $k$ 的无套利取值（或区间）。
\item 若 $k$ 取使某个 $q_s=0$ 或 $q_s<0$ 的值，则分别沿该状态价格的负方向构造组合：当组合成本为零而支付非负且一项为正时为直接套利；当成本为负而支付为零时为融资套利；当成本和支付方向相反时买入低价组合、卖出高价组合。
\item 两个欧式看涨期权的支付向量分别为
\[
Y^{(1)}=(0,0,2)^{\mathsf T},\qquad Y^{(3)}=(0,3/2,1/2)^{\mathsf T}.
\]
完全市场下其价格分别为 $q^{\mathsf T}Y^{(1)}$ 与 $q^{\mathsf T}Y^{(3)}$。
\item 若四状态等概率 $p_s=1/4$，随机贴现因子为 $m_s=q_s/p_s=4q_s$。
\end{enumerate}
\end{solution}
\begin{problem}
\problemno{5.14} 证券市场上有两支股票，股票 A 和股票 B。股票 A 的期初价格为 \num{50} 元，期末两种状态 1 和 2 下的支付分别为 \num{60} 元和 \num{40} 元；股票 B 的期初价格为 \num{40} 元，期末两种状态下的支付分别为 \num{52} 元和 \num{28} 元。
\begin{enumerate}[label=(\alph*)]
    \item 给出价格向量和支付矩阵；
    \item 复制状态或有证券；
    \item 若状态 1 发生的概率为 $80\%$，则等权重组合的期望收益是多少？
\end{enumerate}
\end{problem}
\begin{solution}
\begin{enumerate}[label=(\alph*)]
\item 价格向量为 $S=(50,40)^{\mathsf T}$，支付矩阵为 $X=\begin{bmatrix}60&52\\40&28\end{bmatrix}$。
\item 复制状态 1 证券需解 $60\theta_A+52\theta_B=1$、$40\theta_A+28\theta_B=0$；解得 $\theta_A=-0.07,\theta_B=0.10$，其价格为 $50(-0.07)+40(0.10)=0.5$。复制状态 2 证券解得 $\theta_A=0.13,\theta_B=-0.15$，价格同样为 $0.5$。
\item 等权组合期初价格为 $45$，两状态支付分别为 $56$、$34$。状态 1 概率为 $0.8$ 时，期末期望支付为 $0.8\times56+0.2\times34=51.6$，故期望收益率为 $51.6/45-1=14.67\%$。
\end{enumerate}
\end{solution}
